\scanpage{67}
- \mspage{31} -
of subobjects of $B$ and any $g:G\longrightarrow\mathrm{l.u.b}\,\{A_i\}$ there is a $g':G\longrightarrow A_i$, for some $i$, such that $f_i g'=g$, then $G$ is small. If, in addition, $G$ is projective, then $\varphi_{G,D}$ is a bijection for all directed systems $D$. The dual characterization holds for cosmall objects.

\uline{Proof}: Let $k_{ji}:B_i\longrightarrow B_j$ be any directed system with $k_j:B_j\longrightarrow\operatorname{dir\,lim}B_j$ the induced morphisms. Let $h:G\longrightarrow B_i$ be such that $k_i h:G\longrightarrow\operatorname{dir\,lim}B_j$ is the zero morphism. We must show that $h\in\operatorname{Hom}_{\underline A}(G,B_i)$ represents zero in $\operatorname{dir\,lim}\operatorname{Hom}_{\underline A}(G,B_j)$. Since $k_i h=0$, $k_i h$ factors through $\operatorname{Ker}k_i=\operatorname{dir\,lim}\operatorname{Ker}k_{ji}=\mathrm{l.u.b.}(\operatorname{Ker}k_{ji})$. Since $\{\operatorname{Ker}k_{ji}\}$ is an increasing direct family of subobjects of $A_i$, $k_i h$ factors through some $\operatorname{Ker}k_{ji}$, $j\geq i$. Hence $k_{ji}h=0$ for some $j$. q.e.d.

Suppose now that $G$ is also projective. Given any morphism $h:G\longrightarrow\operatorname{dir\,lim}B_i=\mathrm{l.u.b.}\,(\operatorname{Im}k_i)$, $h$ factors through some $\operatorname{Im}k_i$. Since $B_i\longrightarrow\operatorname{Im}k_i$ is an epimorphism, $h$ lifts to $h':G\longrightarrow B_i$. The class represented by $h'$ in $\operatorname{dir\,lim}\operatorname{Hom}_{\underline A}(G,B_i)$ is mapped by $\varphi_G$ onto $h$. Hence $\varphi_G$ is a bijection.
\mspage{66}
